% =========================================================
%  Capítulo 1 — Introdução
% =========================================================

\chapter{Introdução}
\label{cap:introducao}

O ambiente universitário impõe aos estudantes a necessidade de gerenciar,
simultaneamente, múltiplas disciplinas, avaliações, atividades extracurriculares e
compromissos pessoais. Essa sobrecarga cognitiva frequentemente resulta em baixa
produtividade, procrastinação e abandono dos estudos \cite{[REFERENCIA_1]}.
Ferramentas digitais de organização pessoal têm sido amplamente adotadas como
alternativa para mitigar esses problemas, mas muitas carecem de mecanismos que
promovam o engajamento contínuo do usuário.

A gamificação -- aplicação de elementos e mecânicas de jogos em contextos não
lúdicos -- surge como uma abordagem promissora para aumentar a motivação e a
persistência em plataformas educacionais \cite{[REFERENCIA_2]}. Ao incorporar
sistemas de pontuação, níveis e recompensas ao processo de organização acadêmica,
é possível criar um ciclo de \textit{feedback} positivo que incentiva o estudante
a manter hábitos de estudo regulares.

% -----------------------------------------------------------
\section{Contextualização e Motivação}
\label{sec:contextualizacao}

[Descreva aqui o cenário que motivou a criação do projeto. Inclua dados sobre
a dificuldade dos estudantes universitários em gerir a própria vida acadêmica,
pesquisas sobre procrastinação, evasão escolar, etc.]

% -----------------------------------------------------------
\section{Objetivos}
\label{sec:objetivos}

\subsection{Objetivo Geral}
\label{subsec:objetivo-geral}

Desenvolver uma plataforma \textit{web} gamificada para gestão de estudos
universitários, fornecendo ferramentas de organização acadêmica integradas a
mecânicas de engajamento baseadas em gamificação.

\subsection{Objetivos Específicos}
\label{subsec:objetivos-especificos}

\begin{itemize}
  \item Implementar o módulo de organização acadêmica, contemplando o cadastro e
        gerenciamento de disciplinas, períodos letivos, atividades, avaliações e
        horários de aula;

  \item Projetar e implementar um sistema de gamificação com pontos de experiência
        (XP), níveis progressivos, moedas virtuais e controle de sequência de estudo
        (\textit{streak});

  \item Desenvolver um módulo de sessões de foco (\textit{Focus Sessions}) baseado
        na técnica Pomodoro para recompensar o tempo de estudo dedicado;

  \item Construir o \textit{backend} da aplicação adotando a Arquitetura Hexagonal
        (\textit{Ports \& Adapters}) para garantir o isolamento da lógica de negócio
        da infraestrutura;

  \item Expor os recursos da aplicação por meio de uma API REST segura, documentada
        com Swagger e protegida por autenticação JWT;

  \item Implementar um sistema de notificações para alertar o usuário sobre prazos
        de atividades, limites de faltas e conquistas de gamificação.
\end{itemize}

% -----------------------------------------------------------
\section{Justificativa}
\label{sec:justificativa}

[Explique por que este trabalho é relevante. Relacione a necessidade de ferramentas
de gestão acadêmica com o potencial da gamificação para aumentar o engajamento.
Justifique a escolha das tecnologias (Java/Spring Boot, arquitetura hexagonal).]

% -----------------------------------------------------------
\section{Estrutura do Trabalho}
\label{sec:estrutura}

Este trabalho está organizado em seis capítulos. O Capítulo~\ref{cap:referencial}
apresenta o referencial teórico que fundamenta as escolhas tecnológicas e
metodológicas. O Capítulo~\ref{cap:metodologia} descreve a metodologia de
desenvolvimento adotada e as ferramentas utilizadas. O
Capítulo~\ref{cap:desenvolvimento} detalha a implementação do sistema, cobrindo
arquitetura, modelagem de domínio e principais funcionalidades. O
Capítulo~\ref{cap:resultados} apresenta os resultados obtidos. Por fim, o
Capítulo~\ref{cap:conclusao} conclui o trabalho e aponta direções para trabalhos
futuros.
